\section*{第 \ref{ch:b1:angular-momentum} 章——角动量}

\begin{solution}{exo:b1:angular-momentum:1}
$Fd = 50 \times 0.30 = \qty{15}{N.m}$; $Fd\sin 60^\circ = \qty{13}{N.m}$;
沿手柄方向，作用线通过螺母：$0$。
\end{solution}

\begin{solution}{exo:b1:angular-momentum:2}
$L = mrv = 5.97\times10^{24} \times 1.50\times10^{11} \times 2.98\times10^4 =
\qty{2.7e40}{kg.m^2/s}$; $\dd A/\dd t = L/2m = rv/2 = \qty{2.2e15}{m^2/s}$.
\end{solution}

\begin{solution}{exo:b1:angular-momentum:3}
$r = \ell\sin 30^\circ = \qty{0.50}{m}$, $\omega = \qty{3.4}{rad/s}$:
$L_z/m = r^2\omega = \qty{0.84}{m^2/s}$。张力的作用线
通过支点（对过支点的任一轴力矩为零）；
重力平行于竖直轴（对该轴力矩为零）：$L_z$
守恒。
\end{solution}

\begin{solution}{exo:b1:angular-momentum:4}
$J_1 = 4.0 + 2 \times 3.0 \times 0.64 = \qty{7.8}{kg.m^2}$, $J_2 = 4.0 +
2 \times 3.0 \times 0.04 = \qty{4.2}{kg.m^2}$; $\omega_2 = \omega_1J_1/J_2 =
1.85\omega_1$：\qty{3.7}{}圈每秒。$E_k = L^2/2J$：
比值 $J_1/J_2 = 1.85$；肌肉将手臂向内拉，
克服离心趋势做功。
\end{solution}

\begin{solution}{exo:b1:angular-momentum:5}
$L_\Delta = m\ell^2\dot\theta$；张力：力矩为零（与轴相交）；
重力：$-mg\ell\sin\theta$（力臂 $\ell\sin\theta$，起恢复作用）。
$m\ell^2\ddot\theta = -mg\ell\sin\theta$.
\end{solution}

\begin{solution}{exo:b1:angular-momentum:6}
$r_pv_p = r_av_a$：$v_p/v_a = r_a/r_p = 1.034$；取 $(v_p + v_a)/2 \approx
\qty{29.8}{km/s}$: $v_p = \qty{30.3}{km/s}$, $v_a = \qty{29.3}{km/s}$.
\end{solution}

\begin{solution}{exo:b1:angular-momentum:7}
$L = mrv$ 守恒：$v_2 = v_1r_1/r_2 = \qty{4.0}{m/s}$，$\omega_2 = v_2/r_2
= \qty{16}{rad/s}$（由 \qty{4}{rad/s} 得）。$E_1 = \qty{0.40}{J}$，$E_2 =
\qty{1.6}{J}$：拉力做功 \qty{1.2}{J}。张力沿径向，但
冰球的位移现在具有径向分量（它向内螺旋运动）：功
$\vect T\cdot\dd\vect\ell$ 不再为零。
\end{solution}

\begin{solution}{exo:b1:angular-momentum:8}
对支点的力矩（距支点的距离）：小孩 $30 \times 1.5 =
45$（左侧）；木板的重力作用于其中心，在右侧 $0.5$ 处：$20 \times
0.5 = 10$（右侧）；第二个小孩位于右侧 $x$ 处：$45x$。平衡条件：
$45 = 10 + 45x$，$x = \qty{0.78}{m}$。支点承受的总重力为
$95 \times 9.81 = \qty{932}{N}$.
\end{solution}

\begin{solution}{exo:b1:angular-momentum:9}
两物体：$m_1a = m_1g - T_1$，$m_2a = T_2 - m_2g$；滑轮：$J\dot\omega =
(T_1 - T_2)R$，并取 $a = R\dot\omega$。相加得：$a = (m_1 - m_2)g/(m_1 + m_2 +
J/R^2) = 9.81/(3.0 + 1.0) = \qty{2.5}{m/s^2}$（与 \qty{3.3}{m/s^2} 相比）；
$T_1 = m_1(g - a) = \qty{14.7}{N}$, $T_2 = m_2(g + a) = \qty{12.3}{N}$ ---
张力不相等，因为滑轮需要合力矩。
\end{solution}

\begin{solution}{exo:b1:angular-momentum:10}
$r_1v_1 = r_2v_{\perp 2}$: $v_{\perp 2} = 40 \times 0.5/10 = \qty{2.0}{km/s}$.
面积定律只给出垂直分量；径向分量
需用能量守恒（见\cref{ch:b1:central-forces}）。
\end{solution}

\begin{solution}{exo:b1:angular-momentum:11}
$J\omega$ 守恒，且 $J \propto R^2$：$T' = T(R'/R)^2 = 25 \times 86400
\times (10^4/7\times10^8)^2 = \qty{4.4e-4}{s}$。赤道速率 $2\pi R'/T'
= \qty{1.4e8}{m/s}$，为光速的一半：坍缩无法保持
全部角动量——大部分被抛出（此时相对论开始起作用）。
\end{solution}

\begin{solution}{exo:b1:angular-momentum:12}
$L_z = mr^2\omega$，且 $r$ 变化：$\dd L_z/\dd t = 2mr\dot r\omega \neq 0$。
这必为杆的横向反作用力 $N$ 的力矩：$rN =
2mr\dot r\omega$，因此 $N = 2m\dot r\omega$——杆在珠子
滑出时侧向推动它（即\cref{ch:b1:non-inertial-frames}中的科里奥利项）。
\end{solution}

\begin{solution}{pb:b1:angular-momentum:1}
\textbf{1.} $m\omega^2r = GMm/r^2$: $\omega = \sqrt{GM/r^3} = \sqrt{3.98\times
10^{14}/5.66\times10^{25}} = \qty{2.65e-6}{rad/s}$，周期 $2\pi/\omega =
\qty{2.37e6}{s} = \qty{27.4}{d}$.

\textbf{2.} $L_{\text{轨道}} = 7.35\times10^{22} \times 1.475\times10^{17}
\times 2.65\times10^{-6} = \qty{2.9e34}{kg.m^2/s}$.

\textbf{3.} $v = \sqrt{GM/r}$, $L = mrv = m\sqrt{GMr}$; $\dd L/\dd r =
\tfrac12 m\sqrt{GM/r} = L/2r$.

\textbf{4.} $J = 0.33 \times 5.97\times10^{24} \times 4.06\times10^{13} =
\qty{8.0e37}{kg.m^2}$；$\Omega = \qty{7.29e-5}{rad/s}$；$L_{\text{自转}} =
\qty{5.8e33}{kg.m^2/s}$，为 $L_{\text{轨道}}$ 的五分之一。

\textbf{5.} $J_m \approx 0.4 \times 7.35\times10^{22} \times 3.0\times10^{12} =
\qty{8.9e34}{kg.m^2}$，乘以 $\omega$：\qty{2.4e29}{kg.m^2/s}，为
轨道角动量的 $10^{-5}$，可忽略。

\textbf{6.} 太阳的吸引作用于系统的质心
（一阶近似下对它无力矩）；它对地月这一对天体的潮汐效应
很小：在此所需精度内，总 $L$ 守恒。

\textbf{7.} 较近的隆起部分位于地月连线的前方，
受到月球的拉力，其切向分量与地球的
自转方向相反。

\textbf{8.} 自转上的力矩为负：地球减速，
日长增加。

\textbf{9.} 隆起部分将月球向前拉：轨道上的力矩为正，
$L_{\text{轨道}}$ 增大。

\textbf{10.} $L_{\text{轨道}} \propto \sqrt r$ 增大，故 $r$ 增大；$v \propto
1/\sqrt r$ 减小。

\textbf{11.} 向前的拉力做正功，使月球的能量增加；
轨道越高则速率越小：能量转化为高度（势能），
其增量大于所损失的动能。

\textbf{12.} 地球受潮汐锁定：日长等于月长，
不存在超前的隆起部分，也无力矩。

\textbf{13.} $\Delta\Omega/\Omega = -\qty{2.3e-3}{s}/\qty{86164}{s} =
\num{-2.7e-8}$ 每世纪。

\textbf{14.} $\Delta L_{\text{自旋}} = L_{\text{自旋}}\Delta\Omega/\Omega =
\qty{-1.55e26}{kg.m^2/s}$ 每世纪。

\textbf{15.} $\Delta L_{\text{轨道}} = +\qty{1.55e26}{kg.m^2/s}$ 每世纪。

\textbf{16.} $\Delta r = 2r\,\Delta L/L = 2 \times 3.84\times10^8 \times 1.55
\times10^{26}/2.9\times10^{34} = \qty{4.1}{m}$ 每世纪，\qty{4.1}{cm/yr}：
与实测值 \qty{3.8}{cm/yr} 相近（日长的增长还有
一部分并非潮汐引起）。

\textbf{17.} $\Delta T/T = \tfrac32\Delta r/r = 1.6\times10^{-8}$: $2.37\times
10^6 \times 1.6\times10^{-8} = \qty{0.04}{s}$ 每世纪。

\textbf{18.} $\Delta E_{\text{转动}} = L_{\text{自旋}}\Delta\Omega = 5.8\times10^{33}
\times 7.29\times10^{-5} \times (-2.7\times10^{-8}) = \qty{-1.1e22}{J}$ 每
世纪。

\textbf{19.} $\Delta E_{\text{轨道}} = GMm\Delta r/2r^2 = 2.93\times10^{37}
\times 4.1/(2 \times 1.475\times10^{17}) = \qty{+4.1e20}{J}$ 每世纪。

\textbf{20.} 净值 $-1.1\times10^{22} + 4\times10^{20} \approx \qty{-1.06e22}{J}$
每世纪：$3.4\times10^{12}\ \unit{W}$，由海洋与
固体地球中的潮汐摩擦耗散为热量。

\textbf{21.} 相当于人类功率的五分之一，地热流的
十四分之一——一个虽小却持久的加热器。

\textbf{22.} $L_{\text{tot}} = 3.5\times10^{34}$。当 $r = 5.5\times10^8$ 时：
$\omega = \sqrt{GM/r^3} = \qty{1.55e-6}{rad/s}$, $J\omega = 1.2\times10^{32}$,
$m\sqrt{GMr} = 7.35\times10^{22} \times 4.68\times10^{11} = 3.44\times10^{34}$;
合计 $3.45\times10^{34}$：相符。周期 $2\pi/\omega = \qty{4.1e6}{s}
= \qty{47}{d}$.

\textbf{23.} 太阳的潮汐使地球自转持续慢于月球的
轨道角速度，于是地月之间的力矩会反向；此外，
太阳在此之前早已演化（时间尺度为
数十亿年）。

\textbf{24.} 地球在月球上引起潮汐，使月球的自转受到制动，
直至其一面始终朝向我们；月球较轻，且与较重的
伙伴距离更近，因此月球的锁定要快得多。

\textbf{25.} 守恒的是：总角动量。交换的是：角动量
由地球自转传递到月球轨道，约 \qty{1.6e26}{SI}
每世纪。耗散的是机械能，约几太瓦，
表现为潮汐热。
\end{solution}
